\section{Erd\H{o}s Problem 371}

\paragraph{Problem statement.}
Let $P^{+}(n)$ denote the largest prime factor of $n$. Show that the set of $n$ with $P^{+}(n)<P^{+}(n+1)$ has density $1/2$.

\paragraph{Tools in this round.}
Web browsing: YES. Computation: YES.

\subsection*{1) ROUND-2 OBJECTIVE}
I am not pursuing path (A), because the full natural-density statement is still open. The Round-2 goal is therefore strict advancement: replace Round 1's mere infinitude results by bona fide density theorems, and isolate the exact remaining obstruction between what is known and the desired limit $1/2$.

\subsection*{2) Round-1 FOUNDATION USED}
I rely on two Round-1 facts.
\begin{itemize}
\item For every $n\geq 1$, one has $P^{+}(n)\neq P^{+}(n+1)$, since consecutive integers are coprime.
\item Round 1 gave infinite families in both directions: $P^{+}(2^{m})<P^{+}(2^{m}+1)$ and $P^{+}(2^{m}-1)>P^{+}(2^{m})$.
\end{itemize}
These facts will now be strengthened to genuine density statements.

\subsection*{3) NEW INSIGHT / TOOL (ROUND-2)}
The new ingredients are three published theorems.
\begin{itemize}
\item Ter"av"ainen's window theorem: for any exponents $0<a<b<1$ and $0<c<d<1$, the set of $n$ for which $P^{+}(n)$ and $P^{+}(n+1)$ lie simultaneously in the windows $[n^{a},n^{b}]$ and $[n^{c},n^{d}]$ has positive asymptotic lower density.
\item Ter"av"ainen's logarithmic-density theorem: the set $\{n:P^{+}(n)<P^{+}(n+1)\}$ has logarithmic density $1/2$.
\item Tao--Ter"av"ainen's almost-all-scales theorem, together with the explicit lower-density bound $0.2017$ of L"u--Wang.
\end{itemize}
The first item lets me prove positive lower density for both inequalities directly, which is already a substantial advance over Round 1.

\subsection*{4) ATTACK PLAN (ROUND-2)}
The exact gap left after Round 1 was the absence of any mean-cancellation theorem for the sign sequence
\[
f(n)=\begin{cases}
+1,& P^{+}(n)<P^{+}(n+1),\\
-1,& P^{+}(n)>P^{+}(n+1).
\end{cases}
\]
To prove density $1/2$ one would need
\[
\sum_{n\leq X} f(n)=o(X).
\]
I do not have that. Instead I will do the next best rigorous thing.
\begin{enumerate}
\item Use Ter"av"ainen's window theorem to obtain positive \emph{lower} density for both sides.
\item Record the stronger known theorems: logarithmic density $1/2$, density $1/2$ at almost all scales, and explicit lower density at least $0.2017$.
\item Explain precisely why these do not yet imply the full natural density statement.
\end{enumerate}

\subsection*{5) WORK (ROUND-2)}
Let
\[
S=\{n\geq 1: P^{+}(n)<P^{+}(n+1)\}.
\]
Write
\[
\underline d(A)=\liminf_{X\to\infty} \frac{1}{X}\#(A\cap [1,X])
\]
for asymptotic lower density.

\paragraph{Step 1: a fully proved new theorem -- both sides have positive lower density.}
Ter"av"ainen proved the following.

If $a,b,c,d\in (0,1)$ with $a<b$ and $c<d$, then the set
\[
E_{a,b;c,d}=\{n\geq 1: n^{a}\leq P^{+}(n)\leq n^{b},\; n^{c}\leq P^{+}(n+1)\leq n^{d}\}
\]
has positive asymptotic lower density.

Choose, for example,
\[
a=\frac{1}{10},\qquad b=\frac{1}{5},\qquad c=\frac{4}{5},\qquad d=\frac{9}{10}.
\]
Then $a<b$ and $c<d$, so $E_{a,b;c,d}$ has positive lower density. For every $n\in E_{a,b;c,d}$ we have
\[
P^{+}(n)\leq n^{1/5}<n^{4/5}\leq P^{+}(n+1),
\]
whence $n\in S$. Therefore
\[
\underline d(S)>0.
\]

Now reverse the windows. Let
\[
F=E_{4/5,9/10;1/10,1/5}.
\]
Again Ter"av"ainen's theorem gives $\underline d(F)>0$, and for every $n\in F$ one has
\[
P^{+}(n)\geq n^{4/5}>n^{1/5}\geq P^{+}(n+1),
\]
so $n\notin S$. Hence
\[
\underline d(\mathbf N\setminus S)>0.
\]
This strictly strengthens the Round-1 infinitude result.

\paragraph{Step 2: logarithmic density is exactly $1/2$.}
Ter"av"ainen also proved that the logarithmic density
\[
\delta_{\log}(A)=\lim_{X\to\infty}\frac{1}{\log X}\sum_{\substack{n\leq X\\ n\in A}}\frac{1}{n}
\]
exists for $A=S$ and satisfies
\[
\delta_{\log}(S)=\frac{1}{2}.
\]
So the conjectured symmetry is already known on logarithmic average.

\paragraph{Step 3: at almost all scales, the ordinary density is also $1/2$.}
Tao and Ter"av"ainen proved that there exists an exceptional set $\mathcal X_{0}$ of logarithmic density $0$ such that
\[
\lim_{\substack{X\to\infty\\ X\notin \mathcal X_{0}}}
\frac{1}{X}\#\{n\leq X: n\in S\}=\frac{1}{2}.
\]
Thus the natural-density statement is true for all scales outside a logarithmically negligible exceptional set.

\paragraph{Step 4: the best currently published explicit lower bound.}
L"u and Wang proved
\[
\#\{n<x:P^{+}(n)<P^{+}(n+1)\}>(0.2017-o(1))x,
\]
and the same lower bound for the complementary inequality. Therefore
\[
\underline d(S)\geq 0.2017,
\qquad
\underline d(\mathbf N\setminus S)\geq 0.2017.
\]
Since $S$ and its complement partition the positive integers by the Round-1 no-ties observation, one also gets the corresponding upper bounds
\[
\overline d(S)\leq 0.7983,
\qquad
\overline d(\mathbf N\setminus S)\leq 0.7983.
\]

\subsection*{6) ADVERSARIAL VERIFICATION}
The points most likely to conceal a mistake are the following.
\begin{itemize}
\item \emph{Window inclusion.} In Step 1 the inequalities are strict for $n>1$ because $1/5<4/5$. So every $n$ in the first window set really belongs to $S$, and every $n$ in the reversed window set really belongs to the complement.
\item \emph{Density notions.} Positive lower density, logarithmic density, and ``density at almost all scales'' are three different notions. None of Steps 1--4 silently identifies them.
\item \emph{Why don't Steps 2 and 3 already prove the conjecture?} Because a set of scales of logarithmic density $0$ can still contain subsequences along which the ordinary Ces\`aro averages oscillate. So the Tao--Ter"av"ainen theorem is stronger than anything in Round 1, but weaker than the existence of an ordinary limit along \emph{all} $X$.
\item \emph{Compatibility with Round 1.} The new results strengthen, rather than contradict, Round 1. Infinitude follows immediately from positive lower density.
\end{itemize}

\subsection*{7) FINAL}
\textbf{UNRESOLVED (BUT STRICTLY ADVANCED)}

Round 2 upgrades the Round-1 output substantially:
\begin{itemize}
\item both inequalities now have rigorously proved positive lower density;
\item the set $S$ is known to have logarithmic density $1/2$;
\item the ordinary density equals $1/2$ at almost all scales;
\item the best explicit published lower bound is $0.2017$ for both $S$ and its complement.
\end{itemize}
The exact missing step is now sharply isolated: one must upgrade these partial density results to an ordinary asymptotic density along \emph{all} scales.

\subsection*{8) COMPLETION ESTIMATE (MANDATORY)}
\textbf{COMPLETION: 70\%}

\subsection*{9) REFERENCES}
\begin{itemize}
\item J. Ter"av"ainen, \emph{On binary correlations of multiplicative functions}, Forum Math. Sigma 6 (2018), e10.
\item T. Tao and J. Ter"av"ainen, \emph{The structure of correlations of multiplicative functions at almost all scales, with applications to the Chowla and Elliott conjectures}, Algebra \& Number Theory 13 (2019), no. 9.
\item X. L"u and Z. Wang, \emph{On the largest prime factors of consecutive integers}, Monatsh. Math. 206 (2025), 403--418.
\item T. F. Bloom, \emph{Erd\H{o}s Problem \#371}, accessed 2026-03-07.
\end{itemize}

\subsection*{10) OUTPUT FORMAT}
Erd\H{o}s problem number: $x=371$. This section is written as a drop-in \LaTeX{} fragment and is included in the combined file \texttt{369\_371\_373\_round2\_solutions.tex}.

\bigskip
\hrule
\bigskip

